% LaTeX manuscript skeleton — Gibbs-VLA-Overshoot
\documentclass[11pt]{article}
\usepackage{amsmath,amssymb}
\usepackage{graphicx}
\usepackage{hyperref}

\title{Gibbs Phenomena in Vision-Language-Action Policies:\\
Semantic-to-Physical Impedance and Spectral Signatures}
\author{}
\date{\today}

\begin{document}
\maketitle

\begin{abstract}
We hypothesize that ``jitter'' in Vision-Language-Action (VLA) models is not unstructured noise but a signal-theoretic artifact related to the Gibbs phenomenon: approximating semantic discontinuities with smooth, finite-dimensional controllers yields predictable overshoot and ringing. We outline empirical tests on Language-Table and CALVIN trajectories and relate overshoot ratios to the classical Gibbs constant ($\approx 8.9\%$).
\end{abstract}

\section{Introduction}
% Blind philosopher: high-level logic vs.\ low-level grounding.

\section{Discontinuity in the $L \to A$ map}
% Step-like targets; Dirichlet / partial-sum viewpoint.

\section{Theory}
% sign(x), partial sums S_n, limit ~1.08949 at pi/n.

\section{Methods}
% Trajectories as multivariate time series; FFT / wavelet metrics from repository.

\section{Experiments}
% CALVIN, Language-Table; optional cross-domain weather agent.

\section{Discussion}
% JEPA / manifold flow; discrete-conditioning limitations.

\nocite{*}
\bibliographystyle{plain}
\bibliography{references}

\end{document}
